\section{Conclusion}\label{sec:conclusion}

This work presented HandVerify, a framework for offline handwriting verification based on deep Siamese networks. We systematically evaluated 72 (backbone, loss, dataset-pairing) configurations -- six CNN backbones (ResNet-18/34, EfficientNet-B0/B1, MobileNetV3-Small/Large) trained under three loss paradigms (BCE, Contrastive, Triplet) and assessed on four train$\to$test dataset pairings spanning IAM and RIMES -- establishing baselines and providing practical guidance for verification system design. Each configuration was trained once, with a single random seed and a single writer-disjoint split; no repeated-fold cross-validation was performed.

The principal contributions are as follows: (i) a systematic comparison of three learning paradigms across six CNN backbones, under both same-dataset and cross-dataset evaluation; (ii) the introduction of two training strategies, epoch-wise dynamic negative resampling and selective backbone-layer freezing; (iii) a biometric evaluation using EER, AUC, decidability, and fixed-FAR operating points, with a best same-dataset EER of 3.1\% (IAM$\to$IAM, MobileNetV3-Large + Contrastive) and a best cross-dataset EER of 8.2\% (RIMES$\to$IAM, EfficientNet-B0 + Contrastive); and (iv) a paired statistical analysis (Wilcoxon signed-rank) confirming the consistency of the Contrastive-loss advantage across backbones and splits.

Overall, Contrastive Loss consistently outperformed BCE and Triplet Loss across nearly all backbones, splits, and evaluation regimes, at both the aggregate and best-single-model level, though its advantage was not uniform: gains were largest for backbones where BCE training was unstable (e.g., MobileNetV3-Small) and smallest for ResNet-18, the only backbone with no collapsed BCE configurations. Cross-dataset generalization remained substantially harder than same-dataset evaluation for all losses, particularly under strict low-FAR operating conditions.

\subsection{Future Directions}

Several avenues merit further investigation:

\textbf{Statistical robustness.} Multi-seed training and bootstrap confidence intervals would allow the point estimates reported here to be complemented with proper uncertainty quantification, and would strengthen the independence assumptions underlying the significance analysis.

\textbf{Broader architectural coverage.} Extending the sweep to transformer-based and attention-augmented architectures, and to additional backbone families, could clarify whether the observed loss-function ordering generalizes beyond the six CNN backbones considered here.

\textbf{Advanced negative sampling.} Online hard-negative mining and embedding-distance-based adaptive sampling may further improve Triplet Loss training beyond simple epoch-wise resampling.

\textbf{Fine-grained failure analysis.} A systematic, writer- and pair-level failure analysis across the full test sets -- rather than the single-backbone, top-2-error inspection performed here -- could clarify whether errors concentrate on low-sample-count writers or on genuinely multi-style writers.

\textbf{Explainability.} Visualization techniques such as Grad-CAM could elucidate the regions driving verification decisions, improving interpretability and user trust.

\textbf{Embedding dimensionality analysis.} A systematic study of the embedding dimensionality -- e.g., via PCA-based intrinsic dimensionality estimation or a sweep over the output layer size -- could reveal whether the current embedding size is over- or under-parameterized, and whether a more compact representation would preserve verification performance while reducing storage and matching cost.

\textbf{Multi-modal integration.} Combining offline handwriting with complementary biometric modalities (e.g., signatures, online writing dynamics) may improve overall system robustness.

\textbf{Doddington zoo analysis.} A per-writer categorization into biometric menagerie classes 
(goats, lambs, wolves, and extended classes such as doves, chameleons, 
phantoms and worms~\cite{yager2010biometric}) could reveal 
whether verification errors are broadly distributed across the writer 
population or concentrated in a small subset of problematic identities. 
Such an analysis requires aggregating genuine and impostor scores at the 
individual-writer level rather than pooling them across the test set, which 
was not retained in the current evaluation pipeline; combined with the failure-case 
inspection above, it could clarify whether the false accepts and false rejects observed in Section~5.10 
are symptomatic of a few chronically difficult writers (consistent with the multi-individuality argument 
in Section~5.10) or are spread evenly across the population.

\subsection{Limitations and Confounding Factors}
\label{sec:limitations}

\begin{itemize}
  \item \textbf{Single seed.} Each configuration was trained with a
  single random seed; no per-model bootstrap confidence intervals
  were computed for this revision, so the reported figures are point
  estimates, not interval estimates, of evaluation uncertainty.

  \item \textbf{EER-threshold calibration.} Classification metrics
  (accuracy, precision, recall, F1; Table~\ref{tab:classification-metrics})
  were computed and evaluated on the same set, without an independent
  calibration split, and are therefore mildly optimistic.

  \item \textbf{Single-split same-dataset evaluation.} Same-dataset
  experiments employed a single train/test split rather than
  $k$-fold cross-validation, owing to computational and GPU-hour
  constraints on the Kaggle platform. The reported same-dataset
  performance may therefore be sensitive to the particular writer
  partition, and its variability across alternative splits could not
  be assessed within the present study.

  \item \textbf{Limited data augmentation.} The training pipeline
  employed only RandomResizedCrop, with a relatively narrow scale
  range of $(0.9, 1.1)$, and RandomRotation with a maximum angle of
  $15^\circ$, adopted primarily due to computational constraints.
  Additional geometric and photometric perturbations (e.g., small
  translations, shear, brightness/contrast variation, mild Gaussian
  blur) were not evaluated and remain an open direction.

  \item \textbf{Cross-loss metric comparability.} Raw similarity
  scores are not directly comparable across loss functions in
  absolute terms (BCE yields sigmoid probabilities in $[0,1]$;
  Contrastive/Triplet yield cosine similarities in $[-1,1]$);
  accordingly, only rank-based metrics (AUC, EER, FAR/FRR) and the
  affine-invariant decidability index $d'$ support the comparative
  claims made in this work.

  \item \textbf{Non-Gaussian score distributions.} Non-normality was
  confirmed for all 72 configurations (Section~\ref{sec:distributions});
  $d'$/decidability are therefore used strictly as descriptive
  indices, not as parametric inferential statistics.

  \item \textbf{Non-independence in Wilcoxon testing.} The 24 paired
  observations analyzed in Section~\ref{sec:significance} share
  backbones, data splits, and single-seed training runs; the
  corresponding significance results should be interpreted as
  evidence of directional consistency rather than as 24 statistically
  independent replicates.

  \item \textbf{Granularity of failure-case analysis.} The
  qualitative analysis in Section~\ref{sec:failure-analysis} is
  restricted to distribution- and curve-level patterns for a single
  backbone (ResNet-18), plus a top-2-error pair-level inspection for
  a single (backbone, loss, split) triple; individual false-accept/
  false-reject pairs, writer identifiers, and per-pair scores were
  not systematically retained. A finer-grained, image- and
  writer-level failure analysis is left to future work.

  \item \textbf{Scope of exhaustive validation.} The exhaustive
  impostor-pool validation reported in Section~\ref{sec:exhaustive}
  covers a single (backbone, loss, split) triple (ResNet-18 +
  Contrastive, RIMES$\to$IAM) and does not extend to the remaining
  three splits or to other backbone/loss combinations.

  \item \textbf{Training dynamics.} Per-configuration epoch counts
  and stopping criteria (early stopping vs.\ epoch cap) were not
  available in consolidated form for this revision; raw logs were
  inspected directly only for a subset of BCE/EfficientNet-B1 runs.
\end{itemize}